\documentclass[10pt,a4paper]{amsart}
\usepackage[margin=19mm]{geometry}
\usepackage{amsmath,amssymb,mathtools}
\usepackage[scr=rsfs,cal=euler]{mathalfa}
\usepackage{xfrac}
\usepackage[unicode,hidelinks,bookmarks=false]{hyperref}
\newcommand{\ie}{i.e.\ }
\newcommand{\cf}{cf.\ }
\newcommand{\resp}{respectively\ }
\newcommand{\Subsection}{\subsection}
\providecommand{\nobreakdash}{\nobreak\hbox{-}}
\setlength{\emergencystretch}{2em}
\multlinegap0pt
\usepackage{bm}
\usepackage{bbm}
\usepackage{dsfont}
\usepackage{xspace}
\usepackage{cancel}
\usepackage{mathtools}
\usepackage{amsthm}
\makeatletter
\@ifundefined{gathered}{\newtheorem{gathered}{Gathered}}{}
\@ifundefined{theorem}{\newtheorem{theorem}{Theorem}}{}
\makeatother
\newcommand{\LPT}{{L_2(\mathds{T})}}
\newcommand{\ca}{\mathcal{A}}
\newcommand{\cj}{\mathcal{J}}
\newcommand{\ck}{\mathcal{K}}
\renewcommand{\phi}{\varphi}
\title[On Approximate Representation of Fractional Brownian Motion]{On Approximate Representation of Fractional Brownian Motion}
\author{Konstantin A. Rybakov}
\date{Source: arXiv:2503.04575v3 (CC BY 4.0)}
\begin{document}
\maketitle

\noindent\emph{Source and licence.} On Approximate Representation of Fractional Brownian Motion. Version arXiv:2503.04575v3 (CC BY 4.0), distributed under the Creative Commons Attribution 4.0 International licence, as recorded in the arXiv licence metadata for this exact version and shown on the abstract page https://arxiv.org/abs/2503.04575v3. Full rights notice: licenses/arxiv-2503.04575-CC-BY-4.0.txt.\par\smallskip

\noindent\textit{Editorial scope.} Counted unit: the Theorem whose source label is \texttt{thmMain} in the article (source0.tex lines 293--309, source label \texttt{thmMain}), delivered together with the article's own complete original proof of it (source0.tex lines 311--383, one \texttt{proof} environment of 73 lines). No mathematics is written, repaired or re-derived: statement and proof are the article's own text line for line. Required context retained uncounted: eqDefLeg (lines 116--118); eqDefLegCoef (lines 120--122); eqSpKH (lines 138--140); eqLinearIO (lines 151--153); eqDefFBMFiniteOper (lines 155--162); eqSpFAlphaExplicit (lines 198--202); eqSpFAlphaImplicitAlpha (lines 209--212); eqFracInt (lines 289--291). Every definition, display and label of those lines is retained unchanged. Omitted from this package and not redistributed: 1--115, 119--119, 123--137, 141--150, 154--154, 163--197, 203--208, 213--288, 292--292, 310--310, 384--780. Nothing inside a retained excerpt is omitted or altered: every retained body is delivered exactly as the article's lines are written, so no omission receipt of any kind accompanies this package. Citations: the retained excerpts carry no citation command at all, so no bibliography entry of the article is transcribed and no reference is redistributed. Reference adaptation: none was needed and none was made; every reference inside the delivered unit resolves inside it, so no unresolved reference or citation is produced. Printed slips kept literal: no printed word, symbol, hypothesis or number of the counted statement or of its proof was altered. Layout and class: the article's own class is \texttt{amsart} with packages including amsmath, amssymb, amsfonts, dsfont, cite, graphicx, hyperref; the wrapper is the lane's pinned \texttt{amsart} editorial wrapper at 10pt on A4, which restates the theorem-like environments the retained text needs and adds the article's own macro definitions that the retained text needs (\texttt{\textbackslash LPT}, \texttt{\textbackslash ca}, \texttt{\textbackslash cj}, \texttt{\textbackslash ck}, \texttt{\textbackslash phi}), with extra wrapper packages (bm, bbm, dsfont, xspace, cancel, mathtools). The delivered numbering is the unit's own. Assets: no retained span contains a figure environment, a graphics-inclusion command, a PDF-page inclusion, a table or a TikZ picture; no image file is copied into this package, the package declares no asset dependency and no image file is redistributed with it. Licence: version arXiv:2503.04575v3 of this article is distributed under the Creative Commons Attribution 4.0 International licence, as recorded in the arXiv licence metadata for this exact version and shown on the abstract page https://arxiv.org/abs/2503.04575v3. A full rights notice is delivered at licenses/arxiv-2503.04575-CC-BY-4.0.txt. Characters: every retained excerpt is pure ASCII, so the delivered engine, XeLaTeX with its default Latin Modern text font, reports no missing character.
\begin{equation}\label{eqDefLeg}
  \hat P(i,t) = \sqrt{\frac{2i+1}{T}} \sum\limits_{k=0}^i {l_{ik} \, \frac{t^k}{T^k}}, \ \ \ i = 0,1,2,\dots,
\end{equation}

\begin{equation}\label{eqDefLegCoef}
  l_{ik} = (-1)^{i-k} C^i_{i+k} C^{i-k}_i = (-1)^{i-k} \prod\limits_{m=1}^k \frac{(i-m+1)(i+m)}{m^2}.
\end{equation}

\begin{equation}\label{eqSpKH}
  K_{ij}^H = \int_{\mathds{T}^2} \hat P(i,t) \hat P(j,\tau) k_H(t,\tau) dt d\tau = \int_\mathds{T} \hat P(i,t) \biggl[ \int_0^t k_H(t,\tau) \hat P(j,\tau) d\tau\biggr] dt, \ \ \ i,j = 0,1,2,\dots
\end{equation}

\begin{equation}\label{eqLinearIO}
  \ck_H \phi(t) = \int_0^{t} k_H(t,\tau) \phi(\tau) d\tau, \ \ \ \phi(\cdot) \in \LPT,
\end{equation}

\begin{equation}\label{eqDefFBMFiniteOper}
  \ck_H = \left\{
    \begin{array}{ll}
      a_H \, \cj_{0+}^{2H} \circ \ca_{f_{1/2-H}} \circ \cj_{0+}^{1/2-H} \circ \ca_{f_{H-1/2}} & \text{for} ~ H < 1/2 \\
      a_H \, \cj_{0+}^1 \circ \ca_{f_{H-1/2}} \circ \cj_{0+}^{H-1/2} \circ \ca_{f_{1/2-H}} & \text{for} ~ H > 1/2,
    \end{array}
  \right.
\end{equation}

\begin{equation}\label{eqSpFAlphaExplicit}
  \begin{gathered}
    F_i^\alpha = \int_\mathds{T} {t^\alpha \hat P(i,t) dt} = \sqrt{\frac{2i+1}{T}} \sum\limits_{k=0}^i {l_{ik} \, \frac{T^{\alpha+1}}{\alpha+k+1}} = T^\alpha \sqrt{T} \, \frac{\sqrt{2i+1} \, \alpha^{\underline{i}}}{(\alpha + 1)^{\overline{i+1}}}, \ \ \ i = 0,1,2,\dots,
  \end{gathered}
\end{equation}

\begin{gather}
  F_{i+1}^\alpha = \sqrt{\frac{2i+3}{2i+1}} \, \frac{\alpha-i}{\alpha+i+2} \, F_i^\alpha, \label{eqSpFAlphaImplicit} \\
  F_i^{\alpha+k} = \frac{T^k [(\alpha+1)^{\overline{k}}]^2}{(\alpha-i+1)^{\overline{k}}(\alpha+i+2)^{\overline{k}}} \, F_i^\alpha, \label{eqSpFAlphaImplicitAlpha}
\end{gather}

\begin{equation}\label{eqFracInt}
  \cj_{0+}^\beta t^\alpha = \frac{1}{\Gamma(\beta)} \int_0^t {\frac{\tau^\alpha d\tau}{(t-\tau)^{1-\beta}}} = \frac{\Gamma(\alpha+1)}{\Gamma(\alpha+\beta+1)} \, t^{\alpha+\beta}, \ \ \ \alpha > -1.
\end{equation}

\begin{theorem}\label{thmMain}
Let $K^H$ be the infinite matrix with elements $K_{ij}^H$ defined by Equation~\eqref{eqSpKH}, $H \in (0,1)$. Then, for the $j$th column of the matrix $K^H$, we have
\begin{equation}\label{eqSpOperKHLExplicit}
  K_{*j}^H = a_H \Gamma(3/2-H) \sqrt{\frac{2j+1}{T}} \sum\limits_{k=0}^j \frac{l_{jk} (3/2-H)^{\overline{k}}}{T^k (H+1/2+k) k!} \, F^{H+1/2+k}, \ \ \ j = 0,1,2,\dots,
\end{equation}
where $(3/2-H)^{\overline{k}}$ is the upper factorial, and elements of the infinite column matrix $F^{H+1/2+k}$ are given by Equations~\eqref{eqSpFAlphaExplicit}--\eqref{eqSpFAlphaImplicitAlpha}.

For $H \neq 1/2$, values $K_{ij}^H$ satisfy the following equation:
\begin{equation}\label{eqSpOperKHijExplicit}
  K_{ij}^H = a_H \Gamma(1/2-H) \sqrt{\frac{2j+1}{T}} \, F_i^{H+1/2} \sum\limits_{k=0}^j (-1)^{j-k} \, \frac{1/2-H+k}{H+1/2+k} \, \Pi_k^{(ij)},
\end{equation}
where
\begin{gather*}
  \Pi_0^{(ij)} = 1, \ \ \ \Pi_k^{(ij)} = \frac{(H+1/2+k)^2 (k-1/2-H)}{k^3} \, \frac{j-k+1}{H+1/2-i+k} \, \frac{j+k}{H+1/2+i+k+1} \, \Pi_{k-1}^{(ij)}, \\
  k = 1,2,\dots
\end{gather*}
\end{theorem}

\begin{proof}
Expansion coefficients $K_{ij}^H$ are determined by Equation~\eqref{eqSpKH}. According to Equation~\eqref{eqLinearIO}, it can be rewritten as
\begin{align*}
  K_{ij}^H = \int_\mathds{T} \hat P(i,t) \ck_H \hat P(j,t) dt = \sqrt{\frac{2j+1}{T}} \int_\mathds{T} {\hat P(i,t) \sum\limits_{k=0}^j {l_{jk} \, \frac{\ck_H t^k}{T^k}} \, dt}, \ \ \ i,j = 0,1,2,\dots,
\end{align*}
where the operator $\ck_H$ is defined by Equation~\eqref{eqDefFBMFiniteOper}. Further, we consider three cases.

1. The case $H < 1/2$:
\[
  K_{ij}^H = a_H \sqrt{\frac{2j+1}{T}} \int_\mathds{T} {\hat P(i,t) \sum\limits_{k=0}^j {l_{jk} \, \frac{\cj_{0+}^{2H} \circ \ca_{f_{1/2-H}} \circ \cj_{0+}^{1/2-H} \circ \ca_{f_{H-1/2}} t^k}{T^k}} \, dt}.
\]

Based on Equation~\eqref{eqFracInt}, we have
\begin{align*}
  & \cj_{0+}^{2H} \circ \ca_{f_{1/2-H}} \circ \cj_{0+}^{1/2-H} \circ \ca_{f_{H-1/2}} t^k = \cj_{0+}^{2H} \circ \ca_{f_{1/2-H}} \circ \cj_{0+}^{1/2-H} t^{H-1/2+k} \\
  & \ \ \ = \frac{\Gamma(H+1/2+k)}{\Gamma(k+1)} \, \cj_{0+}^{2H} \circ \ca_{f_{1/2-H}} t^k = \frac{\Gamma(H+1/2+k)}{k!} \, \cj_{0+}^{2H} t^{1/2-H+k} \\
  & \ \ \ = \frac{\Gamma(H+1/2+k) \Gamma(3/2-H+k)}{\Gamma(H+3/2+k) k!} \, t^{H+1/2+k} = \frac{\Gamma(3/2-H+k)}{(H+1/2+k) k!} \, t^{H+1/2+k},
\end{align*}
hence,
\begin{align*}
  K_{ij}^H & = a_H \sqrt{\frac{2j+1}{T}} \sum\limits_{k=0}^j \frac{l_{jk}}{T^k} \, \frac{\Gamma(3/2-H+k)}{(H+1/2+k) k!} \int_\mathds{T} {\hat P(i,t) t^{H+1/2+k} dt} \\
  & = a_H \Gamma(3/2-H) \sqrt{\frac{2j+1}{T}} \sum\limits_{k=0}^j \frac{l_{jk}}{T^k} \, \frac{(1/2-H+1) \ldots (1/2-H+k)}{(H+1/2+k) k!} \, F_i^{H+1/2+k} \\
  & = a_H \Gamma(3/2-H) \sqrt{\frac{2j+1}{T}} \sum\limits_{k=0}^j \frac{l_{jk} (3/2-H)^{\overline{k}}}{T^k (H+1/2+k) k!} \, F_i^{H+1/2+k}, \ \ \ i,j = 0,1,2,\dots,
\end{align*}
where $F_i^{H+1/2+k}$ are expansion coefficients of the power function $f_{H+1/2+k}(t) = t^{H+1/2+k}$ relative to the Legendre polynomials~\eqref{eqDefLeg}. This implies Equation~\eqref{eqSpOperKHLExplicit}.

2. The case $H > 1/2$:
\[
  K_{ij}^H = a_H \sqrt{\frac{2j+1}{T}} \int_\mathds{T} {\hat P(i,t) \sum\limits_{k=0}^j {l_{jk} \, \frac{\cj_{0+}^1 \circ \ca_{f_{H-1/2}} \circ \cj_{0+}^{H-1/2} \circ \ca_{f_{1/2-H}} t^k}{T^k}} \, dt}.
\]

Here, we also use Equation~\eqref{eqFracInt}. Since
\begin{align*}
  & \cj_{0+}^1 \circ \ca_{f_{H-1/2}} \circ \cj_{0+}^{H-1/2} \circ \ca_{f_{1/2-H}} t^k = \cj_{0+}^1 \circ \ca_{f_{H-1/2}} \circ \cj_{0+}^{H-1/2} t^{1/2-H+k} \\
  & \ \ \ = \frac{\Gamma(3/2-H+k)}{\Gamma(k+1)} \, \cj_{0+}^1 \circ \ca_{f_{H-1/2}} t^k = \frac{\Gamma(3/2-H+k)}{k!} \, \cj_{0+}^1 t^{H-1/2+k} = \frac{\Gamma(3/2-H+k)}{(H+1/2+k) k!} \, t^{H+1/2+k},
\end{align*}
we again obtain Equation~\eqref{eqSpOperKHLExplicit}.

3. Case $H = 1/2$:
\[
  K_{ij}^H = \sqrt{\frac{2j+1}{T}} \int_\mathds{T} {\hat P(i,t) \sum\limits_{k=0}^j {l_{jk} \, \frac{\cj_{0+}^1 t^k}{T^k}} \, dt}, \ \ \ \cj_{0+}^1 t^k = \frac{t^{k+1}}{k+1}.
\]

The last equality can be rewritten as
\[
  \cj_{0+}^1 t^k = \frac{\Gamma(3/2-H+k)}{(H+1/2+k) k!} \, t^{H+1/2+k},
\]
since $\Gamma(3/2-H+k) = \Gamma(k+1) = k!$ and $H+1/2+k = k+1$, i.e., Equation~\eqref{eqSpOperKHLExplicit} is valid.

Next, we obtain the equation for elements $K_{ij}^H$ in expanded form. Equation~\eqref{eqSpOperKHLExplicit} implies the following relation:
\[
  K_{ij}^H = a_H \Gamma(3/2-H) \sqrt{\frac{2j+1}{T}} \sum\limits_{k=0}^j \frac{l_{jk} (3/2-H)^{\overline{k}}}{T^k (H+1/2+k) k!} \, F_i^{H+1/2+k}, \ \ \ i,j = 0,1,2,\dots,
\]
in which, according to Equations~\eqref{eqDefLegCoef} and~\eqref{eqSpFAlphaImplicitAlpha}, we have
\begin{align*}
  & \frac{l_{jk} (3/2-H)^{\overline{k}}}{T^k (H+1/2+k) k!} \, F_i^{H+1/2+k} \\
  & \ \ \ = (-1)^{j-k} \frac{(j-k+1) \ldots j(j+1) \ldots (j+k) (1/2-H+1) \ldots (1/2-H+k)}{(k!)^3 (H+1/2+k)} \\
  & \ \ \ \ \ \ {} \times \frac{(H+1/2+1)^2 \dots (H+1/2+k)^2}{(H+1/2-i+1) \dots (H+1/2-i+k) (H+1/2+i+2) \dots (H+1/2+i+k+1)} \, F_i^{H+1/2} \\
  & \ \ \ = \frac{(-1)^{j-k} F_i^{H+1/2}}{H+1/2+k} \prod\limits_{m=1}^k \frac{(H+1/2+m)^2 (1/2-H+m)}{m^3} \, \frac{j-m+1}{H+1/2-i+m} \, \frac{j+m}{H+1/2+i+m+1}.
\end{align*}

Using the Gamma function properties, namely $\Gamma(3/2-H) = (1/2-H) \Gamma(1/2-H)$, and also taking into account the equality
\begin{gather*}
  (1/2-H) \prod\limits_{m=1}^k (1/2-H+m) = (1/2-H) \biggl( \, \prod\limits_{m=1}^{k-1} (1/2-H+m) \biggr) (1/2-H+k) \\
  = (1/2-H+k) \prod\limits_{m=0}^{k-1} (1/2-H+m) = (1/2-H+k) \prod\limits_{m=1}^k (m-1/2-H),
\end{gather*}
we obtain
\begin{align*}
  K_{ij}^H & = a_H \Gamma(1/2-H) \sqrt{\frac{2j+1}{T}} \, F_i^{H+1/2} \sum\limits_{k=0}^j (-1)^{j-k} \, \frac{1/2-H+k}{H+1/2+k} \\
  & \ \ \ {} \times \prod\limits_{m=1}^k \frac{(H+1/2+m)^2 (m-1/2-H)}{m^3} \, \frac{j-m+1}{H+1/2-i+m} \, \frac{j+m}{H+1/2+i+m+1},
\end{align*}
and this is equivalent to Equation~\eqref{eqSpOperKHijExplicit}. The theorem is proved.
\end{proof}

\end{document}
